\section{Auditoría de fuentes: no es una lista de proyectos, sino una clase de diseño de sistemas a escala nacional}

Esta entrevista la impartió Abdullah Alswaha, ministro de Comunicaciones y Tecnologías de la Información de Arabia Saudita. Al verificarlo el 11 de agosto de 2026, el video público ya había pasado a ser private; el repositorio aún conserva 1,266 subtítulos con marca de tiempo y la portada oficial. Por ello, esta sesión reconstruye la argumentación de la clase a partir de los subtítulos y contrasta los mecanismos clave con Saudi Digital Government Authority, KFSHRC, Groq/Aramco Digital y artículos técnicos. En cuanto a las cifras de gigawatt, tamaño de mercado, costos, ahorros y productividad, si ningún documento de primera mano las da con el mismo criterio, este texto las presenta de manera uniforme como \textbf{estimación del ponente}, sin elevar lo dicho en el momento a hechos auditados.

\begin{importantbox}{La cadena causal de tres capas de toda la sesión}
\begin{enumerate}
  \item \textbf{Capa de hardware}: semiconductor, memory, interconnect y cooling determinan la useful compute per watt;
  \item \textbf{Capa de difusión}: la compute capacity solo se convierte en aplicaciones mediante la planificación de recursos, la inference de bajo costo y los workflow de cada sector;
  \item \textbf{Capa de gobernanza}: data, process, human accountability y public value determinan si un agent puede pasar de la demo a los servicios públicos.
\end{enumerate}
\end{importantbox}

\begin{warningbox}{Los tres tipos de evidencia deben separarse}
Las palabras textuales de la clase sirven para reconstruir motivaciones y juicios; los materiales oficiales, para confirmar organizaciones, proyectos o mecanismos vigentes; el análisis de ingeniería, para explicar costos, límites y conclusiones transferibles. Los tres pueden respaldarse entre sí, pero no sustituirse. En particular, la capacidad planeada, el impacto económico y los mercados futuros de los proyectos nacionales solo pueden enunciarse según su nivel de evidencia.
\end{warningbox}

\subsection{Technocrat y signal-to-noise ratio}

La clase describe a Alswaha como \term{technocrat} (tecnócrata): no porque «los ingenieros sean por naturaleza más aptos para gobernar», sino para subrayar que quienes dirigen infraestructura necesitan entender las restricciones de los sistemas, la medición, las fallas y la inversión a largo plazo. El ponente resume su método de juicio con la \term{signal-to-noise ratio} (SNR, relación señal-ruido). Si $P_s$ representa la fuerza de la evidencia relevante para el objetivo y $P_n$ la interferencia y la incertidumbre, puede escribirse
\[
\mathrm{SNR}=\frac{P_s}{P_n}.
\]
Aquí la fórmula es solo una analogía para la toma de decisiones: aumentar la SNR no significa ignorar las objeciones, sino definir primero la función objetivo, la calidad de los datos, la reversibilidad y el costo de fallar, y después extraer señales estables de la propaganda, las fluctuaciones de corto plazo y las cifras no verificadas.

\teachervoice{Alswaha describe su trayectoria profesional con learn--earn--return y considera «volver a Stanford» una oportunidad para transmitir problemas a la siguiente generación. Lo que de verdad vale la pena conservar del tono de la clase no es el relato de éxito personal, sino esto: el talento técnico debe dirigir su capacidad hacia problemas con valor público a largo plazo que, a la vez, puedan verificarse con ingeniería.}

\subsection{Resumen de la sección}

El núcleo de esta sesión no es evaluar la estrategia de un país, sino aprender cómo un sistema de IA a escala nacional se descompone en problemas verificables de hardware, aplicaciones y gobernanza. Todos los casos posteriores vuelven al mismo criterio: cuál es la entrada, dónde está el cuello de botella, cómo se mide el resultado y quién asume la falla.
